% !TeX root = ../main.tex

\chapter{Technische architectuur, alternatieven en lokale opslag}
\label{chap:technische-uitwerking}

\section{Technische architectuur}
\label{chap:architectuur}

\subsection{Electron}
De applicatie is gebouwd met Electron en JavaScript. Electron combineert een Chromium-renderer met
een Node.js-mainproces. Deze keuze maakte het mogelijk om een webgebaseerde interface te bouwen en
tegelijk lokale bestanden en meerdere desktopvensters te beheren. Een bijkomend voordeel was dat dit compatibel was voor beide MacOS en Windows wat nodig was aangezien de app op Windows moest kunnen draaien en de ontwikkeling op MacOS gebeurde.

Het mainproces in \code{src/main/main.js} beheert de levenscyclus van de applicatie, instellingen,
vensters, polling, opslag en IPC-handlers. Een verborgen livevenster laadt de echte timingwebsite.
Een extractiescript wordt in die pagina uitgevoerd om tabelkoppen, cellen en sessie-informatie uit
de gerenderde DOM te lezen.

\subsection{Proces- en modulegrenzen}
De code is bewust in vier zones verdeeld:

\begin{itemize}
  \item \code{src/main}: Electron-specifieke infrastructuur, vensters, updates en bestands-I/O.
  \item \code{src/shared}: pure domeinfuncties voor parsing, analyse, voorspelling, pitstops en
        viewmodellen.
  \item \code{src/renderer}: HTML, CSS en code die reeds berekende data op het scherm zet.
  \item \code{tests}: alle testfiles om domeinfuncties als ook renderer en main functionaliteiten grondig te testen.
\end{itemize}

De scheiding tussen renderer en domeinlogica was een belangrijke verbetering tijdens deze stage.
Vroege versies bevatten berekeningen in de UI-code. Daardoor was het moeilijk om te bewijzen dat een
waarde correct was en konden verschillende schermen dezelfde berekening anders uitvoeren. In de
huidige architectuur leveren modules zoals \code{lapAnalytics.js}, \code{classBattle.js},
\code{lapPrediction.js} en \code{pitstopPlanner.js} volledige resultaten. De renderer beperkt zich
tot presentatie en gebruikersinteractie.

\subsection{Vensterarchitectuur}
Er is één hoofdvenster\ref{fig:versie1} voor de primaire wagen. Wanneer de gebruiker via de setup een tweede of derde
wagen toevoegt, opent voor elke extra wagen een dashboardvenster met dezelfde renderer en een andere
wagenparameter. Alle vensters ontvangen dezelfde collectorstatus. De timingwebsite wordt dus slechts
eenmaal geladen en bevraagd.

Grafieken worden in afzonderlijke vensters geopend.\ref{fig:grafiekvenster1} Ook zij gebruiken de reeds opgeslagen lap
history en starten geen eigen collector. Deze aanpak houdt het hoofdscherm compact en maakt het
mogelijk om analysevensters alleen te openen wanneer ze nodig zijn.

\subsection{Datastroom}
Elke poll doorloopt dezelfde stappen:

\begin{enumerate}
  \item Het extractiescript leest tabellen en sessievelden uit de livepagina.
  \item De parser herkent bruikbare headers en zet elke rij om naar canonieke velden.
  \item De storage-normalizer voegt provider- en sessiecontext toe.
  \item De lap history en actuele rijen worden aan de analysemodules doorgegeven.
  \item Voor elke gevolgde wagen worden dashboardanalyse, voorspelling en pitstopplan berekend.
  \item De volledige collectorstatus wordt opgeslagen en via IPC naar alle vensters gestuurd.
\end{enumerate}

De UI en de berekeningen worden dus gevoed door dezelfde consistente snapshot. Dit voorkomt dat een
grafiek bijvoorbeeld met gegevens van een andere poll rekent dan het hoofdscherm.

\iffalse
\begin{figure}[H]
\centering

\begin{tikzpicture}[
  node distance=0.65cm,
  flowbox/.style={
    draw,
    rounded corners=2pt,
    minimum width=9.5cm,
    minimum height=1.05cm,
    align=center,
    fill=blue!5,
    line width=0.8pt
  },
  outputbox/.style={
    draw,
    rounded corners=2pt,
    minimum width=4.4cm,
    minimum height=1.05cm,
    align=center,
    fill=green!8,
    line width=0.8pt
  },
  arrow/.style={
    -{Latex[length=2.5mm]},
    line width=0.8pt
  }
]

\node[flowbox] (source) {
  \textbf{Live timingwebsite}\\
  RIS Timing of GetRaceResults
};

\node[flowbox, below=of source] (parser) {
  \textbf{Providerspecifieke parser}\\
  Leest de DOM-structuur en herkent kolommen en sessiegegevens
};

\node[flowbox, below=of parser] (normalisation) {
  \textbf{Normalisatie}\\
  Zet providerspecifieke waarden om naar één uniform gegevensmodel
};

\node[flowbox, below=of normalisation] (storage) {
  \textbf{Lokale opslag}\\
  Bewaart snapshots, rondehistoriek, metadata en berekende toestand
};

\node[flowbox, below=of storage] (analytics) {
  \textbf{Analytics en domeinlogica}\\
  Geldigheid, gemiddelden, voorspellingen, gaps en pitstopplanning
};

\node[
  outputbox,
  below=0.8cm of analytics,
  xshift=-4cm
] (dashboard) {
  \textbf{Dashboard en grafieken}\\
  Live informatie voor de race-engineer
};

\node[
  outputbox,
  below=0.8cm of analytics,
  xshift=4cm
] (reports) {
  \textbf{PDF-rapporten}\\
  Analyse per stint en volledige sessie
};

\draw[arrow] (source) -- (parser);
\draw[arrow] (parser) -- (normalisation);
\draw[arrow] (normalisation) -- (storage);
\draw[arrow] (storage) -- (analytics);
\draw[arrow] (analytics) -- (dashboard);
\draw[arrow] (analytics) -- (reports);

\end{tikzpicture}

\caption{Datastroom van de live timingpagina naar het dashboard en de rapportering (verwerken data elk op eigen manier).}
\label{fig:datastroom}
\end{figure}

%% ZET DIT IN DE TEXT
%% Figuur~\ref{fig:datastroom} toont hoe de ruwe timingdata stapsgewijs
%% wordt omgezet naar informatie voor het dashboard en de PDF-rapporten.

\fi

\subsection{Instellingen en configuratie}
Instellingen worden als JSON in de Electron user-datafolder bewaard. Ze omvatten timing-URL,
gevolgde wagens, modus, opslagmap, referentietijden en pitstopregels. Referentietijden zijn
afzonderlijk per sessiemodus opgeslagen, zodat een qualifyingreferentie een racereferentie niet
overschrijft. Circuitprofielen voor Spa, Zolder en Assen bevatten standaardwaarden voor de relevante
pitlengte en FCY-snelheid. In de pitstopsetup kunnen deze waarden per evenement worden overschreven
wanneer het reglement afwijkt.

\section{Afweging van alternatieve oplossingen}
\label{chap:alternatieven}

\subsection{Webapplicatie}
Een centrale webapplicatie zou updates en toegang voor meerdere gebruikers vereenvoudigen, maar vereist extra hosting, authenticatie en een betrouwbare netwerkverbinding. Een lokale browserapplicatie vermijdt die afhankelijkheden, maar biedt minder ondersteuning voor bestandsbeheer en meerdere vensters.

Electron maakte één gedeelde implementatie voor macOS en Windows mogelijk en vereenvoudigde de ontwikkeling en distributie. Een native Windowsapplicatie in C\# zou nog efficiënter kunnen zijn, maar dat was voor dit prototype minder geschikt door de ontwikkeling op verschillende platformen.

\subsection{SQLite}
SQLite was aanvankelijk de voor de hand liggende keuze. Het biedt transacties, indexes en krachtige
queries zonder een afzonderlijke server. Voor de uiteindelijke schaal bleken bestanden echter voordelen
te hebben. Een gebruiker kan een CSV direct openen, JSONL kan append-only worden geschreven en een
beschadigde laatste regel tast niet noodzakelijk de volledige historiek aan. Ook export en debugging
zijn eenvoudiger.
Het nadeel is dat relaties en constraints in de applicatiecode worden beheerd en analyses over meerdere evenementen minder efficiënt zijn. Voor een toekomstig centraal racearchief zou SQLite of PostgreSQL daarom geschikter zijn.

\subsection{Herberekende statistiek}
Een gemiddelde kan efficiënt worden bijgewerkt met een teller en som. Voor een lange race lijkt dit
aantrekkelijk. In dit project zijn ronden echter niet altijd definitief geldig op het eerste moment
dat ze verschijnen. Pitstatus, sectorvlaggen of correcties kunnen de selectie beïnvloeden. Een
gecachet gemiddelde zou dan teruggedraaid moeten worden met exact dezelfde eerdere context.
Daarom wordt de lap history als bron gebruikt en worden statistieken opnieuw uit de gefilterde
relevante ronden berekend. Met enkele duizenden records en pure JavaScriptfuncties blijft dit snel
genoeg. De laatst berekende samenvatting wordt wel opgeslagen voor inspectie. Pas wanneer metingen
aantonen dat herberekenen een werkelijk prestatieprobleem vormt, is een incrementele cache met
invalidatieregels verantwoord.

\subsection{Machine learning model}
Een regressiemodel, random forest of neuraal netwerk zou meerdere features kunnen combineren:
piloot, bandenslijtage, brandstofniveau, verkeer, temperatuur en weersomstandigheden. De
timingpagina levert veel van die inputs niet. De beschikbare trainingsdata is bovendien beperkt en
bevat veranderende circuits en omstandigheden. Een complex model zou daardoor snel overfitten en
moeilijk uitlegbaar zijn.

\section{Dataverwerking en lokale opslag}
\label{chap:data}

De betrouwbaarheid van het dashboard hangt af van een consistente verwerking en opslag van de ontvangen timingdata. Gegevens van verschillende timingproviders worden daarom eerst omgezet naar één intern formaat en vervolgens lokaal bewaard. Zo kunnen analyses steeds dezelfde datastructuur gebruiken en blijft de racehistoriek ook beschikbaar na een onderbreking of herstart.

\subsection{Uniforme parser}
De verschillende timingproviders gebruiken verschillende labels en representaties. De parser begint daarom met het
opschonen van witruimte en het omzetten van headers naar een standaardvorm. Varianten worden gekoppeld aan interne
sleutels voor positie, startnummer, klasse, PIC, team, piloot, aantal ronden, gap, interval,
sectoren, laatste tijd, beste tijd en pitinformatie.

Tijdswaarden worden zo vroeg mogelijk naar milliseconden omgezet. De tekst \code{2:05.073} wordt
intern bijvoorbeeld \code{125073 milliseconden}. Formatteringsfuncties zetten dit pas bij presentatie
opnieuw om. Dit maakt optellen, gemiddeldes nemen, vergelijken en testen eenvoudiger.

\subsection{Vlagstatus en pitfasen}
Voor elke sector wordt de vlagstatus vastgelegd op het moment dat de sector beschikbaar wordt. Dit
detail is nodig voor het scenario waarin sector 1 groen werd gereden en tijdens sector 2 een FCY
begint. De volledige ronde is dan niet representatief voor een rondetijdgemiddelde, maar sector 1
kan nog geldig zijn voor het gemiddelde en de beste waarde van sector 1.
De analyselaag annoteert ook pitfasen. Veranderingen in pitteller en pitstatus helpen inlaps en
outlaps identificeren. Deze ronden blijven deel van de historische werkelijkheid, maar worden niet
gebruikt voor bijvoorbeeld gemiddeldes te berekenen. Dit voorkomt dat een trage pitronde het gemiddelde van een
piloot onnodig verslechtert.

\subsection{Bestandsformaten}
De belangrijkste bestanden zijn:

\begin{itemize}
  \item \code{lap_history.jsonl}: append-only historiek met 1 JSON-object per opgeslagen ronde
  \item \code{lap_history.csv}: dezelfde rondehistoriek in tabelvorm voor manuele inspectie of spreadsheetanalyse
  \item \code{latest_live_rows.json} en \code{latest_live_rows.csv}: laatste volledige timingsnapshot
  \item \code{session_metadata.json}: sessiecontext, gevolgde wagen(s), timingprovider en instellingen voor pitstops
  \item \code{analytics_summary.json}: laatst berekende statistieken voor het dashboard en de grafieken
  \item \code{lap_prediction.json} en \code{lap_prediction_car-<nummer>.json}: actuele rondevoorspelling per gevolgde wagen
  \item \code{pitstop_plan.json} en \code{pitstop_plan_car-<nummer>.json}: actuele pitstopplanning per gevolgde wagen
  \item \code{gap_state.json} en \code{gap_history.jsonl}: laatst bevestigde gapinformatie en gap-historiek
  \item \code{stint_state.json}: actuele stintinformatie per wagen en piloot
  \item \code{track_condition_state.json} en \code{track_condition_events.jsonl}: huidige baanconditie en wijzigingen doorheen de sessie
\end{itemize}

JSONL is geschikt voor rondes omdat een nieuwe ronde kan worden toegevoegd zonder het volledige
bestand te herschrijven. Tegelijk blijft elke regel leesbaar en herstelbaar. Samenvattingsbestanden
worden wel overschreven: ze vertegenwoordigen de actuele afgeleide toestand en kunnen indien nodig
opnieuw uit de lap history worden berekend.
